% Chapter 2
\chapter{کاربرد‌های دوقلوی دیجیتال}
دوقلوی دیجیتال کاربرد‌های زیادی در زمینه‌های مختلف دارد که در این بخش به توضیح مختصری از هرکدام پرداخته ‌شده است.\cite{b14}
\section{دوقلوی دیجیتال پزشکی‌ و سلامت}

با نگاهی گذرا به اطراف،‌ می‌توانیم نمونه‌های فراوانی از دستگاه‌ها و حسگر‌هایی را مشاهده کرد که در حال نظارت و بررسی وضعیت سلامت و فیزیولوژیکی انسان‌ها هستند. این دستگاه‌ها با پایش و بررسی دائمی یا موقت وضعیت جسمانی افراد، آن‌ را گزارش کرده و تحلیل می‌کند. برای نمونه می‌توان به ساعت‌های هوشمند، دستگاه‌های تصویربرداری پزشکی، نرم‌افزار‌‌های دیجیتال سلامت، دماسنج اشاره کرد. هرکدام از این دستگاه‌ها داده‌های فراوانی که مربوط به بخش‌های مختلف بدن و سلامت افراد هست تولید می‌کنند. این داده‌ها ممکن است با هم در ارتباط باشند و به بررسی و نظارت و پایش\LTRfootnote{Monitoring} آن‌ها به صورت یکجا نیاز باشد. همچنین علاوه بر مشاهده نیاز به تحلیل داده‌ها و ارسال اطلاعات مورد نیاز به پزشک متخصص است. این کار می‌تواند به کمک دوقلوی دیجیتال به همراه دریافت اطلاعات از دستگاه‌های اینترنت اشیاء انجام شود. استفاده از دوقلوی دیجیتال در اهدای عضو، آموزش جراحی و کاهش خطر سایر روش‌های پزشکی، رو به توسعه است.


\section{شهرهای هوشمند}

دوقلوی دیجیتال شهر یک مفهوم جدید است که مدلی از برنامه‌ریزی و ساخت و ساز شهری برای توسعه پایدار آینده ارائه می‌دهد که به طور موثری نوآوری‌ها در فناوری دیجیتال را با مکانیسم‌های عملیاتی شهری ترکیب و مسیری امکان‌پذیر برای ارتقاء شهری فراهم می‌کند. 
با نقشه‌برداری‌های دقیق از شهر، ارتباط واقعیت و دنیای مجازی با یکدیگر و بازخورد هوشمند نمونه‌ فیزیکی و دیجیتالی شهر، فعالیت‌ها و پدیده‌های شهری به صورت کارآمد‌تر پربازده‌تر انجام می‌شود. همچنین خدمات روزمره برای شهروندان راحت‌تر و فراگیر‌تر می‌شود و به ایجاد یک محیط پایدار کمک می‌کند.

در حالی ک شهر‌ها در دنیا، به سرعت در مسیر رشد هستند، در این مسیر دچار چالش‌هایی نیز می‌شوند. رکود اقتصادی، افزایش آلودگی‌های هوا و زیست محیطی، تراکم ترافیک شهری، بلایای طبیعی مانند طوفان، سیل و زلزله، می‌توانند نمونه‌هایی از چالش‌هایی باشد که حتی شهر‌های مدرن نیز با آن‌ها دست‌و پنجه نرم می‌کنند. دوقلوی دیجیتال شهر از منظر‌های مختلف، می‌تواند یک راه‌حل بالقوه برای سیاست‌گذاران و تصمیم‌گیران در جهان باشد که به‌وسیله آن، در جهت بهبود مدیریت شهری و ایجاد یک بوم‌سازگان\RTLfootnote{Ecosystem} شهری مساعد برای صنایع و مردم حرکت کنند.

استفاده از دوقلوی دیجیتال در شهر‌ها می‌تواند میزان صرفه‌جویی هزینه‌ها را افزایش دهد و بازگشت سرمایه را به سرعت، با خود به همراه داشته‌ باشد. پس می‌توان به سادگی پیش‌بینی کرد که در آینده برای هر شهر، مدیریت بحران‌های موجود و تحلیل آن به وسیله دوقلوی دیجیتال، یک امر ضروری است.

\section{صنعت ساخت و تولید}
فناوری دوقلوی دیجیتال، به طور گسترده در صنعت تولید استفاده می‌شود. دوقلو‌ی دیجیتال در صنعت تولید، فرایند‌ها و رخداد‌ها، پدیده‌ها و اتفاق‌ها را واضح و شفاف می‌کند. صنعت تولید و ساخت‌وساز را می‌توان مهد دوقلوی دیجیتال دانست زیرا اولین استفاده‌ها و نیازمندی‌ها در راستای صنعت۴\RTLfootnote{Industry 4} برای دوقلوی دیجیتال، در این صنعت ایجاد و انجام شده‌ است.

دوقلوی دیجیتال در صنعت تولید، می‌تواند با پویش یکپارچه و بررسی دائمی اجزاء و آلات موجود در صنعت، از پیشامد‌های ناگوار پیشرو جلوگیری کند و باعث کاهش چشم‌گیر هزینه‌های مختلف شود. استفاده از فناوری‌های مربوط به اینترنت اشیاء در این صنعت بسیار رایج است. بنابر این رفتن به سمت دوقلوی دیجیتال در آن بیشتر از زمینه‌های دیگر مورد انتظار است. همچنین خودکار‌سازی در این صنعت از اهمیت به خصوصی برخوردار است که با استفاده از دوقلوی دیجیتال می‌تواند فرایند‌های خودکار‌سازی کامل‌تر و به‌صرفه تر شود.


\section{خودرو}
دوقلوی دیجیتال می‌تواند مدل مجازی یک وسیله نقلیه فیزیکی را ایجاد کند. دوقلوی دیجیتال وسایل نقلیه می‌تواند با جمع‌آوری داده‌های مختلف از سنسور‌های متنوع موجود در آن، یک کپی مجازی از آن را ایجاد کند. با داشتن دوقلوی دیجیتال خودرو، می‌توان پیش‌بینی‌های لازم، جهت اتخاذ بهترین تصمیم به صورت خودکار یا توسط راننده، انجام شود. 

خودرو‌های خودران نیز در همین زمینه قابل بیان هستند. این خودرو‌ها برای اینکه به مرحله خودران بودن برسند نیاز به ساعت‌ها آموزش توسط داده‌های واقعی و شبیه‌سازی‌های موجود هستند. ایجاد یک دوقلوی مجازی برای خودرو و آموزش آن توسط داده‌هایی که خود، می‌تواند دوقلوی مجازی از سایر خودرو‌ها و جاده‌ها باشد، این عمل را ممکن می‌کند. با این شرایط، می‌توان خودرو‌های خودران را ابتدا در دنیای مجازی ایجاد و رشد داد و پس از آنکه آزمون خود را در دنیای مجازی پس دادند و تا حد امکان بدون خطا عمل کردند، در دنیای واقعی شروع به کار کنند.

دوقلوی دیجیتال در خودرو‌ها را می‌توان به سطح شهر و خیابان‌ها آورد. به این صورت که یک دوقلوی مجازی از خودرو‌های موجود در خیابان به طوری که با هم در ارتباط باشند، ایجاد کرد. به عبارتی دیگر، یک دوقلوی مجازی چند عاملی که در آن عامل‌ها، خودرو‌ها هستند که با هم تبادل اطلاعات می‌کنند و بر اساس آن،‌ تصمیم‌گیری کرده و عملی را انجام می‌دهند. مثلا می‌توانند فاصله مناسب را از یکدیگره به صورت ایمن، حفظ کنند و یا مسیر‌های دارای ترافیک را انتخاب نکنند و به صورتی هرکدام مسیر خود را برگزیند، که کمترین ترافیک در سطح شهر ایجاد شود.

می‌توان دوقلوی دیجیتال خودرو را برای هر کدام از قطعات مکانیکی خودرو به صورت دقیق ایجاد کرد. می‌دانیم که قطعات خودرو در کنار هم می‌توانند یک وسیله نقلیه کامل را ایجاد کنند. پس یک دوقلوی مجازی چند عاملی می‌توان برای هر خودرو داشت که هر عامل آن، یک قطعه مهم و حیاتی در آن، در نظر گرفته می‌شود و به صورت مستقل، بررسی و پایش شرایط لحظه‌ای آن مهم است و در صورت پیش‌بینی ایجاد مشکل یا خرابی در آن، می‌توان در زمان مناسب، آن را تعمیر یا تعویض کرد پیش از آنکه هزینه کلان یا غیر قابل جبرانی ایجاد کند.

\section{انسان دیجیتال}
ایجاد انسان مجازی یا دیجیتال و دوقلوی مجازی انسان\LTRfootnote{Human Digital Twin} 
یکی دیگر از مواردی است که در سال‌های اخیر مورد توجه قرار گرفته است. این نوع از انسان، نه تنها به صورت آواتارهایی هوشمند بجای منشی، در خدمت شرکت‌ها هستند بلکه با رشد هوش مصنوعی، تا حدی توانسته‌اند کاملا حسی مشابه حس ارتباط با انسان واقعی را القاء کنند.

انسان دیجیتال به طور گسترده به عنوان یک برنامه رایانه‌ای تعریف می‌شود که یک ظاهر مانند انسان را ارائه می‌دهد. این انسان مجازی، در واقع گرافیک رایانه‌ای، بینایی رایانه، گفتار هوشمند، پردازش زبان طبیعی و سایر فناوری‌ها را یکپارچه می‌کند. می‌توان از آن برای تولید محتوای دیجیتال و تعامل انسان و رایانه برای کمک به بهبود کارآیی تولید محتوا و تجربه کاربر استفاده کرد. در واقع دوقلوی دیجیتال انسان، یک انسان در دنیای مجازی و غیر فیزیکی است که به صورت یک محصول جامع، با ابزار‌های مختلف مانند یادگیری عمیق،‌ پردازش گفتار و نوشتار، توانایی تعامل با انسان و غیر انسان ایجاد می‌شود و کاربرد‌های فراوانی در زمینه‌های مختلف می‌تواند داشته باشد.

\section{جمعیت مجازی}
با ایجاد دوقلوی مجازی فرد به طوری که کاملا مانند او فکر کند، حرکت کند، تصمیم بگیرد و حتی واکنش‌های فیزیکی و شیمیایی بدن مشابه داشته باشد و همچنین جسم آن، به صورت مشابه با نمونه واقعی در برابر دارو‌ها و خوراک‌ها عمل کند، می‌توان به سمت ایجاد یک جامعه مجازی حرکت کرد.

بنابر این تکثیر دوقلوی دیجیتال بدن انسان می‌تواند در بخش‌های مختلف از اهمیت بالایی برخوردار باشد. برای مثال اگر فرایند درمان یک بیماری خاص در نظر گرفته شود، می‌توان فرایند درمان را در سه مرحله بررسی کرد:
\begin{enumerate}
	\item 
ساخت نسخه‌های نامحدود از مدل‌های شبکه از انسان و همه عوامل مولکولی به صورت کاملا مشابه با انسان و بیماری مربوطه.
	\item 
درمان محاسباتی در دوقلوی مجازی با هزاران دارو به منظور شناسایی بهترین دارو برای بیماری مورد نظر.
	\item 
درمان بیمار با داروی مشخص شده پس از تست‌های بسیار در نمونه‌های دیجیتالی.
\end{enumerate}

\section{صنایع هوافضا}
ایجاد دوقلوی دیجیتال برای صنایع هوافضا به ویژه هواپیما و اجزای پیچیده آن را می‌تواند یکی از پر اهمیت‌ترین کاربرد‌های دوقلوی دیجیتال دانست زیرا که هزینه‌های خرابی و تعمیر در این صنعت بسیار بالا است و این درحالی است که به کمک دوقلوی دیجیتال و ایجاد دوقلوی دیجیتال قطعات صنعت هوافضا، می‌توان بسیاری از خرابی‌ها و مشکل‌های آن‌ها را پیش‌بینی کرده و از خرابی آن‌ها جلوگیری کرد.

همچنین استفاده از دوقلوی دیجیتال به عنوان شبیه‌ساز پرواز برای هواپیما‌ها و سایر اقلام پرنده، می‌تواند در پیش‌بینی حالت‌های مختلف پرواز موثر باشد. همچنین می‌توان از ایجاد دوقلوی مجازی آن‌ها در آموزش خلبان‌ها استفاده شود و می‌توان شبیه‌سازی شرایط سخت و غیرممکن پیش از اینکه در واقعیت ایجاد شود، را ساخت.

\section{صنایع نفت، گاز و پتروشیمی}
در این صنعت یا به طور کلی، صنعت هیدروکربن، طیف وسیعی از فناوری‌های دیجیتال را برای بهبود بهره‌وری، کارآیی و ایمنی در نظر می‌گیرد و در عین حال، هزینه‌های سرمایه و عملیاتی، خطرات بهداشتی و زیست محیطی به میزان قابل توجهی کاهش می‌یابد. این منابع بسیار ارزشمند هستند و در جهت جلوگیری از هدر رفت آن در فرایند‌های مختلف تولیدی و انتقال،‌ باید اقدام شود و دوقلوی مجازی پایگاه‌های استخراج و راه‌‌های انتقال می‌تواند گزینه بسیار مناسبی برای این کار باشد.

همچنین قرارگیری ایستگاه‌های انتقال در مکان‌های دور از دسترس و سختی بررسی آن و هزینه بالای مربوط به نگهداری و مراقبت از سامانه‌های مورد نیاز در این صنعت، استفاده از دوقلوی مجازی را بیشتر مفید می‌سازد.

